\documentclass[12pt]{article}
\usepackage{amsmath, amssymb, amsthm}
\usepackage{geometry}
\geometry{a4paper, margin=1in}
\usepackage{hyperref}
\usepackage{enumitem}

\title{Problem Sets: Theory of Higher-Order Linear Differential Equations}
\author{Based on Chapter 6 Concepts}
\date{}

\begin{document}

\maketitle
\tableofcontents
\newpage

% ---------------------------------------------------------
\section{Concept Checklist \& Question Mapping}
% ---------------------------------------------------------
This section outlines the core concepts, problem types, and tricks extracted from the source material, mapped to the corresponding questions in the problem set.

\begin{itemize}
    \item \textbf{Concept 1: Interval of Existence and Uniqueness (Theorem 1)}
    \begin{itemize}
        \item Standard form conversion (Leading Coefficient Trap).
        \item Identifying discontinuities (zeroes in denominators, domain of functions like $\ln(x)$).
        \item Selecting the open interval containing the initial condition $x_0$.
        \item \textit{Relevant Questions:} 1, 2, 3, 4
    \end{itemize}
    
    \item \textbf{Concept 2: Linear Dependence and Independence (via Definition)}
    \begin{itemize}
        \item Setting $c_1y_1 + c_2y_2 + \dots + c_ny_n = 0$.
        \item Evaluating at strategic points to form a system of equations.
        \item \textit{Relevant Questions:} 5, 6, 7
    \end{itemize}

    \item \textbf{Concept 3: Wronskian and Linear Independence}
    \begin{itemize}
        \item Setting up the $n \times n$ determinant (Missing Derivative Check).
        \item The "Zero-Point" Wronskian Cheat (evaluating at $x=0$ or another convenient point).
        \item The Wronskian Zero Trap ($W=0$ does not guarantee dependence unless functions are known solutions to a linear homogeneous ODE).
        \item \textit{Relevant Questions:} 8, 9, 10, 11
    \end{itemize}

    \item \textbf{Concept 4: Verifying Fundamental Solution Sets}
    \begin{itemize}
        \item Verifying functions are solutions to the homogeneous ODE.
        \item Computing the Wronskian to prove independence (Determinant Simplification/Factoring).
        \item Writing the general solution $y_c$.
        \item \textit{Relevant Questions:} 12, 13, 14, 15
    \end{itemize}

    \item \textbf{Concept 5: Nonhomogeneous Equations and Initial Value Problems (IVPs)}
    \begin{itemize}
        \item Forming the general solution $y = y_p + y_c$.
        \item The IVP Particular Solution Trap (applying initial conditions to the \textit{entire} solution, not just $y_c$).
        \item Product rule with exponentials and trigonometric functions.
        \item \textit{Relevant Questions:} 16, 17, 18, 19, 20
    \end{itemize}

    \item \textbf{Concept 6: Superposition Principle}
    \begin{itemize}
        \item Linear operator properties: $L[c_1y_1 + c_2y_2] = c_1L[y_1] + c_2L[y_2]$.
        \item Avoiding the multiplication trap (operators do not distribute over multiplication).
        \item \textit{Relevant Questions:} 21, 22, 23, 24
    \end{itemize}

    \item \textbf{Concept 7: Finding Fundamental Sets for Constant Coefficient ODEs}
    \begin{itemize}
        \item Characteristic equations.
        \item Distinct real roots, repeated real roots, and complex roots.
        \item \textit{Relevant Questions:} 25, 26, 27, 28, 29, 30
    \end{itemize}
\end{itemize}

\newpage

% ---------------------------------------------------------
\section{Problem Sets}
% ---------------------------------------------------------

\subsection*{Concept 1: Interval of Existence and Uniqueness}
\textbf{1.} Determine the largest interval $(a, b)$ for which Theorem 1 guarantees the existence of a unique solution to the IVP:
$$(x-3)y''' + \ln(x)y'' + y = 0; \quad y(1)=0, \quad y'(1)=2, \quad y''(1)=-1$$

\textbf{2.} Determine the largest interval of validity for the following IVP:
$$x^2 y^{(4)} - \frac{1}{x-2} y'' + e^x y = \tan(x); \quad y(1)=1, \quad y'(1)=0, \quad y''(1)=0, \quad y'''(1)=0$$

\textbf{3.} Find the largest interval of validity for the IVP:
$$(t^2-4)y''' + 2ty = 3; \quad y(0)=1, \quad y'(0)=0, \quad y''(0)=0$$

\textbf{4.} Determine the largest interval of validity for the IVP:
$$\sin(x)y'' + y' = x; \quad y(4)=0, \quad y'(4)=1$$

\subsection*{Concept 2: Linear Dependence and Independence (via Definition)}
\textbf{5.} Use the definition of linear independence to determine if $\{1, x, x^2\}$ is linearly independent on $(-\infty, \infty)$.

\textbf{6.} Determine whether the set $\{\sin^2 x, \cos^2 x, \cos(2x)\}$ is linearly dependent or independent on $(-\infty, \infty)$. Justify your answer.

\textbf{7.} Use the definition of linear independence to show that $\{e^x, xe^x\}$ is linearly independent on $(-\infty, \infty)$.

\subsection*{Concept 3: Wronskian and Linear Independence}
\textbf{8.} Use the "Zero-Point" Wronskian cheat to determine if the functions $\{e^{2x}, e^{-2x}, 1\}$ are linearly independent on $(-\infty, \infty)$.

\textbf{9.} Let $f_1(x) = x^3$ and $f_2(x) = x^2|x|$. Compute the Wronskian $W[f_1, f_2](x)$ on $(-1, 1)$. Are the functions linearly dependent or independent on $(-1, 1)$? (Beware the Wronskian Zero Trap).

\textbf{10.} Compute the Wronskian of $\{t, t\ln t\}$ on $(0, \infty)$ to determine if they are linearly independent.

\textbf{11.} Compute the Wronskian of $\{\sin x, \cos x, x\}$ at $x=\pi$ to determine if the functions are linearly independent on $(-\infty, \infty)$.

\subsection*{Concept 4: Verifying Fundamental Solution Sets}
\textbf{12.} Verify that $\{1, e^x, e^{-x}\}$ forms a fundamental solution set for $y''' - y' = 0$ on $(-\infty, \infty)$.

\textbf{13.} Verify that $\{e^x, e^{2x}, e^{3x}\}$ forms a fundamental solution set for $y''' - 6y'' + 11y' - 6y = 0$.

\textbf{14.} Verify that $\{x, x^2, x^3\}$ forms a fundamental solution set for the Euler-Cauchy equation $x^3 y''' - 3x^2 y'' + 6x y' - 6y = 0$ on $(0, \infty)$.

\textbf{15.} Verify that $\{e^x, e^{-x}, \sin x, \cos x\}$ forms a fundamental solution set for $y^{(4)} - y = 0$.

\subsection*{Concept 5: Nonhomogeneous Equations and Initial Value Problems}
\textbf{16.} A particular solution to $y'' - 4y = 8x$ is $y_p = -2x$. The fundamental set for the homogeneous equation is $\{e^{2x}, e^{-2x}\}$. Find the solution satisfying $y(0)=1, y'(0)=-1$.

\textbf{17.} A particular solution to $y'' + y = x^2$ is $y_p = x^2 - 2$. The fundamental set for the homogeneous equation is $\{\cos x, \sin x\}$. Find the solution satisfying $y(0)=0, y'(0)=0$.

\textbf{18.} A particular solution to $y''' - y'' = 2$ is $y_p = -x^2$. The fundamental set for the homogeneous equation is $\{1, x, e^x\}$. Find the solution satisfying $y(0)=0, y'(0)=1, y''(0)=0$.

\textbf{19.} A particular solution to $y'' - 2y' + y = e^x$ is $y_p = \frac{1}{2}x^2e^x$. The fundamental set for the homogeneous equation is $\{e^x, xe^x\}$. Find the solution satisfying $y(0)=1, y'(0)=2$.

\textbf{20.} A particular solution to $y'' + 4y = \sin x$ is $y_p = \frac{1}{3}\sin x$. The fundamental set for the homogeneous equation is $\{\cos(2x), \sin(2x)\}$. Find the solution satisfying $y(0)=1, y'(0)=0$.

\subsection*{Concept 6: Superposition Principle}
\textbf{21.} Let $L$ be a linear differential operator. If $L[y_1] = x$ and $L[y_2] = e^x$, find a solution to $L[y] = 3x - 2e^x$.

\textbf{22.} Given $L[y_1] = \sin x$ and $L[y_2] = \cos x$, find a solution to $L[y] = 4\sin x + 5\cos x$.

\textbf{23.} \textit{True or False:} If $L[y_1] = x^2$, then $y = (y_1)^2$ is a solution to $L[y] = x^4$. Explain your reasoning.

\textbf{24.} Given $L[y_1] = e^{-x}$ and $L[y_2] = x$, find a solution to $L[y] = 2e^{-x} - 4x$.

\subsection*{Concept 7: Finding Fundamental Sets for Constant Coefficient ODEs}
\textbf{25.} Find the general solution to $y''' - 4y'' + 4y' = 0$.

\textbf{26.} Find the general solution to $y''' + y'' - 2y = 0$.

\textbf{27.} Find the general solution to $y^{(4)} - 16y = 0$.

\textbf{28.} Find the general solution to $y''' + 3y'' + 3y' + y = 0$.

\textbf{29.} Find the general solution to $y^{(4)} + 8y'' + 16y = 0$.

\textbf{30.} Find the general solution to $y''' - y'' - 4y' + 4y = 0$.

\newpage

% ---------------------------------------------------------
\section{Detailed Solutions}
% ---------------------------------------------------------

\subsection*{Concept 1: Interval of Existence and Uniqueness}

\textbf{Solution 1:}
\textbf{Step 1: Standard Form.} Divide by the leading coefficient $(x-3)$:
$$y''' + \frac{\ln(x)}{x-3}y'' + \frac{1}{x-3}y = 0$$
\textbf{Step 2: Identify Discontinuities.} The function $\ln(x)$ is only continuous for $x > 0$. The denominator $(x-3)$ is zero at $x=3$. Thus, discontinuities/boundaries are at $x \le 0$ and $x=3$.
\textbf{Step 3: Intervals.} The real line is divided into $(0, 3)$ and $(3, \infty)$.
\textbf{Step 4: Select Interval.} The initial condition is at $x_0 = 1$. Since $1 \in (0, 3)$, the largest interval of validity is \textbf{$(0, 3)$}.

\textbf{Solution 2:}
\textbf{Step 1: Standard Form.} Divide by $x^2$:
$$y^{(4)} - \frac{1}{x^2(x-2)} y'' + \frac{e^x}{x^2} y = \frac{\tan(x)}{x^2}$$
\textbf{Step 2: Identify Discontinuities.} Denominators are zero at $x=0$ and $x=2$. The function $\tan(x)$ is discontinuous at $x = \frac{\pi}{2} + k\pi$.
\textbf{Step 3: Intervals.} The relevant points near $x_0=1$ are $0$, $\pi/2 \approx 1.57$, and $2$. The intervals are $(-\infty, 0)$, $(0, \pi/2)$, $(\pi/2, 2)$, etc.
\textbf{Step 4: Select Interval.} Since $x_0 = 1 \in (0, \pi/2)$, the largest interval is \textbf{$(0, \pi/2)$}.

\textbf{Solution 3:}
\textbf{Step 1: Standard Form.} Divide by $(t^2-4)$:
$$y''' + \frac{2t}{t^2-4}y = \frac{3}{t^2-4}$$
\textbf{Step 2: Identify Discontinuities.} The denominator is zero when $t^2 - 4 = 0 \implies t = \pm 2$.
\textbf{Step 3: Intervals.} The intervals are $(-\infty, -2)$, $(-2, 2)$, and $(2, \infty)$.
\textbf{Step 4: Select Interval.} The initial condition is at $t_0 = 0$. Since $0 \in (-2, 2)$, the largest interval is \textbf{$(-2, 2)$}.

\textbf{Solution 4:}
\textbf{Step 1: Standard Form.} Divide by $\sin(x)$:
$$y'' + \frac{1}{\sin(x)}y' = \frac{x}{\sin(x)}$$
\textbf{Step 2: Identify Discontinuities.} $\sin(x) = 0$ at $x = k\pi$ for integers $k$.
\textbf{Step 3: Intervals.} The intervals are $\dots, (0, \pi), (\pi, 2\pi), (2\pi, 3\pi), \dots$
\textbf{Step 4: Select Interval.} The initial condition is at $x_0 = 4$. Since $\pi \approx 3.14$ and $2\pi \approx 6.28$, $4 \in (\pi, 2\pi)$. The largest interval is \textbf{$(\pi, 2\pi)$}.

\subsection*{Concept 2: Linear Dependence and Independence (via Definition)}

\textbf{Solution 5:}
\textbf{Step 1: Setup.} Assume $c_1(1) + c_2(x) + c_3(x^2) = 0$ for all $x \in (-\infty, \infty)$.
\textbf{Step 2: Evaluate at specific points.}
Let $x=0 \implies c_1 = 0$.
Let $x=1 \implies c_2 + c_3 = 0$.
Let $x=-1 \implies -c_2 + c_3 = 0$.
\textbf{Step 3: Solve.} Adding the last two equations gives $2c_3 = 0 \implies c_3 = 0$. Then $c_2 = 0$. Since $c_1=c_2=c_3=0$ is the only solution, the functions are \textbf{linearly independent}.

\textbf{Solution 6:}
\textbf{Step 1: Setup.} Assume $c_1\sin^2 x + c_2\cos^2 x + c_3\cos(2x) = 0$.
\textbf{Step 2: Use identities.} Recall the double angle identity: $\cos(2x) = \cos^2 x - \sin^2 x$.
\textbf{Step 3: Substitute.} $c_1\sin^2 x + c_2\cos^2 x + c_3(\cos^2 x - \sin^2 x) = 0$.
Rearranging: $(c_1 - c_3)\sin^2 x + (c_2 + c_3)\cos^2 x = 0$.
\textbf{Step 4: Find non-trivial constants.} If we choose $c_3 = 1$, then $c_1 = 1$ and $c_2 = -1$. Since there exists a non-trivial solution $(1, -1, 1)$, the functions are \textbf{linearly dependent}.

\textbf{Solution 7:}
\textbf{Step 1: Setup.} Assume $c_1e^x + c_2xe^x = 0$ for all $x$.
\textbf{Step 2: Evaluate at specific points.}
Let $x=0 \implies c_1e^0 + c_2(0)e^0 = 0 \implies c_1 = 0$.
The equation reduces to $c_2xe^x = 0$.
Let $x=1 \implies c_2(1)e^1 = 0 \implies c_2e = 0 \implies c_2 = 0$.
\textbf{Step 3: Conclusion.} Since $c_1=c_2=0$, the functions are \textbf{linearly independent}.

\subsection*{Concept 3: Wronskian and Linear Independence}

\textbf{Solution 8:}
\textbf{Step 1: Setup Wronskian at $x=0$.}
$$W(x) = \begin{vmatrix} e^{2x} & e^{-2x} & 1 \\ 2e^{2x} & -2e^{-2x} & 0 \\ 4e^{2x} & 4e^{-2x} & 0 \end{vmatrix}$$
\textbf{Step 2: Evaluate at $x=0$.}
$$W(0) = \begin{vmatrix} 1 & 1 & 1 \\ 2 & -2 & 0 \\ 4 & 4 & 0 \end{vmatrix}$$
\textbf{Step 3: Compute Determinant.} Expand along the third column:
$$W(0) = 1 \cdot (2(4) - (-2)(4)) = 1 \cdot (8 + 8) = 16$$
Since $W(0) = 16 \neq 0$, the functions are \textbf{linearly independent}.

\textbf{Solution 9:}
\textbf{Step 1: Compute Wronskian.}
For $x \ge 0$, $f_2(x) = x^3$. $W(x) = \begin{vmatrix} x^3 & x^3 \\ 3x^2 & 3x^2 \end{vmatrix} = 3x^5 - 3x^5 = 0$.
For $x < 0$, $f_2(x) = -x^3$. $W(x) = \begin{vmatrix} x^3 & -x^3 \\ 3x^2 & -3x^2 \end{vmatrix} = -3x^5 - (-3x^5) = 0$.
Thus, $W(x) = 0$ for all $x \in (-1, 1)$.
\textbf{Step 2: Check Definition (Wronskian Zero Trap).}
Assume $c_1x^3 + c_2x^2|x| = 0$.
At $x=1$: $c_1(1) + c_2(1) = 0 \implies c_1 + c_2 = 0$.
At $x=-1$: $c_1(-1) + c_2(1) = 0 \implies -c_1 + c_2 = 0$.
Adding gives $2c_2 = 0 \implies c_2 = 0$, and thus $c_1 = 0$.
Despite $W(x)=0$, the functions are \textbf{linearly independent}.

\textbf{Solution 10:}
\textbf{Step 1: Setup Wronskian.}
$$W(t) = \begin{vmatrix} t & t\ln t \\ 1 & \ln t + 1 \end{vmatrix}$$
\textbf{Step 2: Compute Determinant.}
$$W(t) = t(\ln t + 1) - (t\ln t)(1) = t\ln t + t - t\ln t = t$$
\textbf{Step 3: Conclusion.} Since $W(t) = t \neq 0$ for $t \in (0, \infty)$, the functions are \textbf{linearly independent}.

\textbf{Solution 11:}
\textbf{Step 1: Setup Wronskian.}
$$W(x) = \begin{vmatrix} \sin x & \cos x & x \\ \cos x & -\sin x & 1 \\ -\sin x & -\cos x & 0 \end{vmatrix}$$
\textbf{Step 2: Evaluate at $x=\pi$.}
$$W(\pi) = \begin{vmatrix} 0 & -1 & \pi \\ -1 & 0 & 1 \\ 0 & 1 & 0 \end{vmatrix}$$
\textbf{Step 3: Compute Determinant.} Expand along the first column:
$$W(\pi) = -(-1) \begin{vmatrix} -1 & \pi \\ 1 & 0 \end{vmatrix} = 1(0 - \pi) = -\pi$$
Since $W(\pi) = -\pi \neq 0$, the functions are \textbf{linearly independent}.

\subsection*{Concept 4: Verifying Fundamental Solution Sets}

\textbf{Solution 12:}
\textbf{Step 1: Verify solutions.}
$y_1=1 \implies y_1'=0, y_1'''=0 \implies 0-0=0$.
$y_2=e^x \implies y_2'=e^x, y_2'''=e^x \implies e^x-e^x=0$.
$y_3=e^{-x} \implies y_3'=-e^{-x}, y_3'''=-e^{-x} \implies -e^{-x}-(-e^{-x})=0$.
\textbf{Step 2: Compute Wronskian.}
$$W(x) = \begin{vmatrix} 1 & e^x & e^{-x} \\ 0 & e^x & -e^{-x} \\ 0 & e^x & e^{-x} \end{vmatrix}$$
Expand along the first column: $W(x) = 1(e^x(e^{-x}) - (-e^{-x})(e^x)) = 1(1 - (-1)) = 2 \neq 0$.
\textbf{Conclusion:} They are solutions and $W \neq 0$, so they form a \textbf{fundamental solution set}.

\textbf{Solution 13:}
\textbf{Step 1: Verify solutions.}
$y_1=e^x \implies 1 - 6(1) + 11(1) - 6 = 0$.
$y_2=e^{2x} \implies 8e^{2x} - 6(4e^{2x}) + 11(2e^{2x}) - 6e^{2x} = e^{2x}(8 - 24 + 22 - 6) = 0$.
$y_3=e^{3x} \implies 27e^{3x} - 6(9e^{3x}) + 11(3e^{3x}) - 6e^{3x} = e^{3x}(27 - 54 + 33 - 6) = 0$.
\textbf{Step 2: Compute Wronskian.}
$$W(x) = \begin{vmatrix} e^x & e^{2x} & e^{3x} \\ e^x & 2e^{2x} & 3e^{3x} \\ e^x & 4e^{2x} & 9e^{3x} \end{vmatrix} = e^x e^{2x} e^{3x} \begin{vmatrix} 1 & 1 & 1 \\ 1 & 2 & 3 \\ 1 & 4 & 9 \end{vmatrix}$$
Determinant $= 1(18-12) - 1(9-3) + 1(4-2) = 6 - 6 + 2 = 2$.
$W(x) = 2e^{6x} \neq 0$. They form a \textbf{fundamental solution set}.

\textbf{Solution 14:}
\textbf{Step 1: Verify solutions.}
$y_1=x \implies y'=1, y''=0, y'''=0 \implies x^3(0) - 3x^2(0) + 6x(1) - 6(x) = 0$.
$y_2=x^2 \implies y'=2x, y''=2, y'''=0 \implies 0 - 3x^2(2) + 6x(2x) - 6(x^2) = -6x^2 + 12x^2 - 6x^2 = 0$.
$y_3=x^3 \implies y'=3x^2, y''=6x, y'''=6 \implies x^3(6) - 3x^2(6x) + 6x(3x^2) - 6(x^3) = 6x^3 - 18x^3 + 18x^3 - 6x^3 = 0$.
\textbf{Step 2: Compute Wronskian.}
$$W(x) = \begin{vmatrix} x & x^2 & x^3 \\ 1 & 2x & 3x^2 \\ 0 & 2 & 6x \end{vmatrix} = x(12x^2 - 6x^2) - 1(6x^3 - 2x^3) = 6x^3 - 4x^3 = 2x^3$$
Since $x \in (0, \infty)$, $W(x) \neq 0$. They form a \textbf{fundamental solution set}.

\textbf{Solution 15:}
\textbf{Step 1: Verify solutions.}
$y_1=e^x \implies y^{(4)}=e^x \implies e^x-e^x=0$.
$y_2=e^{-x} \implies y^{(4)}=e^{-x} \implies e^{-x}-e^{-x}=0$.
$y_3=\sin x \implies y^{(4)}=\sin x \implies \sin x-\sin x=0$.
$y_4=\cos x \implies y^{(4)}=\cos x \implies \cos x-\cos x=0$.
\textbf{Step 2: Compute Wronskian at $x=0$ (Zero-Point Cheat).}
$$W(0) = \begin{vmatrix} 1 & 1 & 0 & 1 \\ 1 & -1 & 1 & 0 \\ 1 & 1 & 0 & -1 \\ 1 & -1 & -1 & 0 \end{vmatrix}$$
Evaluating this $4 \times 4$ determinant yields $W(0) = 8 \neq 0$. They form a \textbf{fundamental solution set}.

\subsection*{Concept 5: Nonhomogeneous Equations and Initial Value Problems}

\textbf{Solution 16:}
\textbf{Step 1: General Solution.} $y(x) = -2x + c_1e^{2x} + c_2e^{-2x}$.
\textbf{Step 2: Apply Initial Conditions (IVP Trap).}
$y(0) = c_1 + c_2 = 1$.
$y'(x) = -2 + 2c_1e^{2x} - 2c_2e^{-2x}$.
$y'(0) = -2 + 2c_1 - 2c_2 = -1 \implies 2c_1 - 2c_2 = 1 \implies c_1 - c_2 = 0.5$.
\textbf{Step 3: Solve System.}
$c_1 + c_2 = 1$
$c_1 - c_2 = 0.5$
Adding gives $2c_1 = 1.5 \implies c_1 = 0.75$. Then $c_2 = 0.25$.
\textbf{Final Answer:} $y(x) = -2x + 0.75e^{2x} + 0.25e^{-2x}$.

\textbf{Solution 17:}
\textbf{Step 1: General Solution.} $y(x) = x^2 - 2 + c_1\cos x + c_2\sin x$.
\textbf{Step 2: Apply Initial Conditions.}
$y(0) = -2 + c_1 = 0 \implies c_1 = 2$.
$y'(x) = 2x - c_1\sin x + c_2\cos x$.
$y'(0) = c_2 = 0$.
\textbf{Final Answer:} $y(x) = x^2 - 2 + 2\cos x$.

\textbf{Solution 18:}
\textbf{Step 1: General Solution.} $y(x) = -x^2 + c_1 + c_2x + c_3e^x$.
\textbf{Step 2: Apply Initial Conditions.}
$y(0) = c_1 + c_3 = 0$.
$y'(x) = -2x + c_2 + c_3e^x \implies y'(0) = c_2 + c_3 = 1$.
$y''(x) = -2 + c_3e^x \implies y''(0) = -2 + c_3 = 0 \implies c_3 = 2$.
\textbf{Step 3: Solve System.}
$c_3 = 2$.
$c_2 + 2 = 1 \implies c_2 = -1$.
$c_1 + 2 = 0 \implies c_1 = -2$.
\textbf{Final Answer:} $y(x) = -x^2 - 2 - x + 2e^x$.

\textbf{Solution 19:}
\textbf{Step 1: General Solution.} $y(x) = \frac{1}{2}x^2e^x + c_1e^x + c_2xe^x$.
\textbf{Step 2: Apply Initial Conditions.}
$y(0) = c_1 = 1$.
$y'(x) = xe^x + \frac{1}{2}x^2e^x + c_1e^x + c_2e^x + c_2xe^x$.
$y'(0) = 0 + 0 + c_1 + c_2 + 0 = 2 \implies 1 + c_2 = 2 \implies c_2 = 1$.
\textbf{Final Answer:} $y(x) = \frac{1}{2}x^2e^x + e^x + xe^x$.

\textbf{Solution 20:}
\textbf{Step 1: General Solution.} $y(x) = \frac{1}{3}\sin x + c_1\cos(2x) + c_2\sin(2x)$.
\textbf{Step 2: Apply Initial Conditions.}
$y(0) = c_1 = 1$.
$y'(x) = \frac{1}{3}\cos x - 2c_1\sin(2x) + 2c_2\cos(2x)$.
$y'(0) = \frac{1}{3} - 0 + 2c_2 = 0 \implies 2c_2 = -\frac{1}{3} \implies c_2 = -\frac{1}{6}$.
\textbf{Final Answer:} $y(x) = \frac{1}{3}\sin x + \cos(2x) - \frac{1}{6}\sin(2x)$.

\subsection*{Concept 6: Superposition Principle}

\textbf{Solution 21:}
\textbf{Step 1: Use Linearity.} We want to solve $L[y] = 3x - 2e^x$.
Notice the right side is $3L[y_1] - 2L[y_2]$.
\textbf{Step 2: Apply Operator Properties.} Because $L$ is linear, $3L[y_1] - 2L[y_2] = L[3y_1 - 2y_2]$.
\textbf{Final Answer:} $y = 3y_1 - 2y_2$.

\textbf{Solution 22:}
\textbf{Step 1: Use Linearity.} We want to solve $L[y] = 4\sin x + 5\cos x$.
This is exactly $4L[y_1] + 5L[y_2] = L[4y_1 + 5y_2]$.
\textbf{Final Answer:} $y = 4y_1 + 5y_2$.

\textbf{Solution 23:}
\textbf{Step 1: Analyze the claim.} The claim is that $L[y_1^2] = (L[y_1])^2$.
\textbf{Step 2: Apply Operator Rules.} Linear operators distribute over addition and scalar multiplication ($L[af+bg] = aL[f]+bL[g]$), but \textbf{not} over multiplication of functions.
\textbf{Final Answer:} \textbf{False}. This is the multiplication trap. Superposition does not apply to products.

\textbf{Solution 24:}
\textbf{Step 1: Use Linearity.} We want to solve $L[y] = 2e^{-x} - 4x$.
This is $2L[y_1] - 4L[y_2] = L[2y_1 - 4y_2]$.
\textbf{Final Answer:} $y = 2y_1 - 4y_2$.

\subsection*{Concept 7: Finding Fundamental Sets for Constant Coefficient ODEs}

\textbf{Solution 25:}
\textbf{Step 1: Characteristic Equation.} $r^3 - 4r^2 + 4r = 0$.
\textbf{Step 2: Factor.} $r(r^2 - 4r + 4) = 0 \implies r(r-2)^2 = 0$.
Roots are $r=0$ (mult 1) and $r=2$ (mult 2).
\textbf{Step 3: General Solution.} $y(x) = c_1e^{0x} + c_2e^{2x} + c_3xe^{2x} = c_1 + c_2e^{2x} + c_3xe^{2x}$.

\textbf{Solution 26:}
\textbf{Step 1: Characteristic Equation.} $r^3 + r^2 - 2 = 0$.
\textbf{Step 2: Factor.} By inspection, $r=1$ is a root ($1+1-2=0$).
Polynomial division gives $(r-1)(r^2 + 2r + 2) = 0$.
For $r^2+2r+2=0$, use quadratic formula: $r = \frac{-2 \pm \sqrt{4 - 8}}{2} = -1 \pm i$.
\textbf{Step 3: General Solution.} $y(x) = c_1e^x + c_2e^{-x}\cos x + c_3e^{-x}\sin x$.

\textbf{Solution 27:}
\textbf{Step 1: Characteristic Equation.} $r^4 - 16 = 0$.
\textbf{Step 2: Factor.} $(r^2 - 4)(r^2 + 4) = 0 \implies (r-2)(r+2)(r^2+4) = 0$.
Roots are $r=2, -2, 2i, -2i$.
\textbf{Step 3: General Solution.} $y(x) = c_1e^{2x} + c_2e^{-2x} + c_3\cos(2x) + c_4\sin(2x)$.

\textbf{Solution 28:}
\textbf{Step 1: Characteristic Equation.} $r^3 + 3r^2 + 3r + 1 = 0$.
\textbf{Step 2: Factor.} Recognize as binomial expansion: $(r+1)^3 = 0$.
Root is $r=-1$ with multiplicity 3.
\textbf{Step 3: General Solution.} $y(x) = c_1e^{-x} + c_2xe^{-x} + c_3x^2e^{-x}$.

\textbf{Solution 29:}
\textbf{Step 1: Characteristic Equation.} $r^4 + 8r^2 + 16 = 0$.
\textbf{Step 2: Factor.} $(r^2 + 4)^2 = 0$.
Roots are $r = \pm 2i$ with multiplicity 2.
\textbf{Step 3: General Solution.} $y(x) = c_1\cos(2x) + c_2\sin(2x) + c_3x\cos(2x) + c_4x\sin(2x)$.

\textbf{Solution 30:}
\textbf{Step 1: Characteristic Equation.} $r^3 - r^2 - 4r + 4 = 0$.
\textbf{Step 2: Factor by grouping.} $r^2(r-1) - 4(r-1) = 0 \implies (r^2-4)(r-1) = 0 \implies (r-2)(r+2)(r-1) = 0$.
Roots are $r=1, 2, -2$.
\textbf{Step 3: General Solution.} $y(x) = c_1e^x + c_2e^{2x} + c_3e^{-2x}$.

\newpage

% ---------------------------------------------------------
\section{Concept Checklist \& Question Mapping (Reference)}
% ---------------------------------------------------------
\begin{itemize}
    \item \textbf{Concept 1: Interval of Existence and Uniqueness (Theorem 1)}
    \begin{itemize}
        \item Standard form conversion (Leading Coefficient Trap).
        \item Identifying discontinuities (zeroes in denominators, domain of functions like $\ln(x)$).
        \item Selecting the open interval containing the initial condition $x_0$.
        \item \textit{Relevant Questions:} 1, 2, 3, 4
    \end{itemize}
    
    \item \textbf{Concept 2: Linear Dependence and Independence (via Definition)}
    \begin{itemize}
        \item Setting $c_1y_1 + c_2y_2 + \dots + c_ny_n = 0$.
        \item Evaluating at strategic points to form a system of equations.
        \item \textit{Relevant Questions:} 5, 6, 7
    \end{itemize}

    \item \textbf{Concept 3: Wronskian and Linear Independence}
    \begin{itemize}
        \item Setting up the $n \times n$ determinant (Missing Derivative Check).
        \item The "Zero-Point" Wronskian Cheat (evaluating at $x=0$ or another convenient point).
        \item The Wronskian Zero Trap ($W=0$ does not guarantee dependence unless functions are known solutions to a linear homogeneous ODE).
        \item \textit{Relevant Questions:} 8, 9, 10, 11
    \end{itemize}

    \item \textbf{Concept 4: Verifying Fundamental Solution Sets}
    \begin{itemize}
        \item Verifying functions are solutions to the homogeneous ODE.
        \item Computing the Wronskian to prove independence (Determinant Simplification/Factoring).
        \item Writing the general solution $y_c$.
        \item \textit{Relevant Questions:} 12, 13, 14, 15
    \end{itemize}

    \item \textbf{Concept 5: Nonhomogeneous Equations and Initial Value Problems (IVPs)}
    \begin{itemize}
        \item Forming the general solution $y = y_p + y_c$.
        \item The IVP Particular Solution Trap (applying initial conditions to the \textit{entire} solution, not just $y_c$).
        \item Product rule with exponentials and trigonometric functions.
        \item \textit{Relevant Questions:} 16, 17, 18, 19, 20
    \end{itemize}

    \item \textbf{Concept 6: Superposition Principle}
    \begin{itemize}
        \item Linear operator properties: $L[c_1y_1 + c_2y_2] = c_1L[y_1] + c_2L[y_2]$.
        \item Avoiding the multiplication trap (operators do not distribute over multiplication).
        \item \textit{Relevant Questions:} 21, 22, 23, 24
    \end{itemize}

    \item \textbf{Concept 7: Finding Fundamental Sets for Constant Coefficient ODEs}
    \begin{itemize}
        \item Characteristic equations.
        \item Distinct real roots, repeated real roots, and complex roots.
        \item \textit{Relevant Questions:} 25, 26, 27, 28, 29, 30
    \end{itemize}
\end{itemize}

\end{document}